\subsection{\bt{Publication and amortization}{发布与摊薄}}
\label{sec:sharing-eval}
\bi{\emph{What admission refuses (Q6).} In the runs of Section~\ref{sec:scenarios}, \PubJudged\ learning trajectories were judged successful; admission published \PubPromoted\ of them as new entries, \PubCorroborated\ corroborated an existing entry instead of adding a copy, and \PubRefused\ were refused. Refusals are rare, but each would have reached every consumer. Over both end-to-end studies, \PubAllJudged\ judged-successful trajectories from three LLMs, the canonical replay refused \PubReplayRefused\ candidates into which extraction had introduced an error, and the provenance check refused \PubChainRefused\ whose recorded derivation does not reproduce the answer. In the worst case, the extractor kept the current-version filter of electronics store revenue but dropped the category filter: the agent had answered \PubExAnswer, and the canonical SQL that later agents would have received returns \PubExCanon, \PubExRatio\ times as much. The producer's own SQL was correct; only checking what consumers will read catches the error.}
{\emph{准入拒绝了什么（Q6）。}在第~\ref{sec:scenarios} 节的运行中，\PubJudged\ 条学习轨迹被判定成功；准入把其中 \PubPromoted\ 条发布为新条目，\PubCorroborated\ 条作为已有条目的佐证而不是再加一份副本，拒绝了 \PubRefused\ 条。拒绝很少，但每一次被放行都会传到每个使用者。在两轮端到端实验中（3 个大模型、\PubAllJudged\ 条判题成功的轨迹），规范 SQL 重放拒绝了 \PubReplayRefused\ 个在提取时被引入错误的候选，溯源检查拒绝了 \PubChainRefused\ 个记录的推导无法复现答案的候选。最严重的一例中，提取器保留了电子品类门店营业额的当前版本过滤，却丢掉了类别过滤：智能体答的是 \PubExAnswer，而之后的智能体将拿到的规范 SQL 返回 \PubExCanon，是它的 \PubExRatio\ 倍。生产者自己的 SQL 是对的；只有检查使用者将要读到的内容，才能发现这个错误。}

\bi{\emph{When sharing pays off.} Building the memory costs, per run, the ten learning sessions (\AmLearnTokK\,k tokens, \AmLearnS\,s), extraction (\AmExtractTokK\,k tokens, \AmExtractS\,s), and admission (\AmAdmitS\,s), in all \AmProdTokK\,k tokens and \AmProdS\,s, counting the learning sessions in full although they also answer their own questions. A held-out task then costs a consumer \AmHoldTokCondK\,k tokens and \AmHoldSCond\,s, against \AmHoldTokAloneK\,k tokens and \AmHoldSAlone\,s when it explores alone, so the memory repays its construction after \AmBreakS\ consumer tasks in time and \AmBreakTok\ in tokens; counted per correct answer (\AmTokPerCorrectCondK\ against \AmTokPerCorrectAloneK\,k tokens), it does so after \AmBreakCorrect\ correct answers. Maintenance is paid once per change for all consumers: it takes \AmMaintBenignS\,s per benign change and \AmMaintRepairS\,s per change whose definitions are repaired. Consumers that verify validity themselves pay per task instead: the self-verifying example-retrieval agent spends \AmVerifyExtraS\,s and \AmVerifyExtraTokK\,k tokens more per task than plain retrieval on unchanged data and benign updates, about \AmTasksPerChange\ tasks per change in our runs, and is still less accurate under breaking changes (Section~\ref{sec:scenarios}).}
{\emph{共享何时回本。}每次运行中，建立记忆的代价是 10 个学习会话（\AmLearnTokK\,k token，\AmLearnS\ 秒）、提取（\AmExtractTokK\,k token，\AmExtractS\ 秒）和准入（\AmAdmitS\ 秒），合计 \AmProdTokK\,k token、\AmProdS\ 秒；学习会话本身也回答了各自的问题，这里仍全额计入。此后每道留出题，使用者花 \AmHoldTokCondK\,k token、\AmHoldSCond\ 秒，独自探索时为 \AmHoldTokAloneK\,k token、\AmHoldSAlone\ 秒，因此按时间计记忆在 \AmBreakS\ 道使用者任务后回本，按 token 计在 \AmBreakTok\ 道后回本；按每个正确答案计（\AmTokPerCorrectCondK\ 对 \AmTokPerCorrectAloneK\,k token），在 \AmBreakCorrect\ 个正确答案后回本。维护每次变化只为所有使用者付一次：每次正常变化 \AmMaintBenignS\ 秒，定义被修复的每次变化 \AmMaintRepairS\ 秒。自己核验有效性的使用者则按任务付费：在数据不变和正常更新时，自行核验的示例检索智能体每题比普通检索多花 \AmVerifyExtraS\ 秒和 \AmVerifyExtraTokK\,k token，我们的运行中每次变化约有 \AmTasksPerChange\ 道任务，而且在破坏性变化下它仍然不如维护准确（第~\ref{sec:scenarios} 节）。}

\bi{\emph{When invalidation returns the work to consumers.} The balance reverses where no repair is possible. Under the duplicate load, the date-key format change, and the backup copy, invalidated definitions send consumers back to the data with a notice that names the failed condition. They then spend \AmInvalTokCondK\,k tokens and \AmInvalSCond\,s per task, against \AmInvalTokAloneK\,k tokens and \AmInvalSAlone\,s for agents that explore alone, and answer \AmInvalAccCond\% against \AmInvalAccAlone\%; tasks on valid or repaired definitions cost \AmValidTokCondK\,k tokens. The fixes these consumers find are paid for in every task; admitting them as candidate repairs would let the layer pay for each fix once.}
{\emph{失效何时把工作交还给使用者。}在无法修复的地方，收支反转。在重复加载、日期键格式变化和备份副本下，失效的定义把使用者送回数据，并附上指明失败条件的通知。此时使用者每题花 \AmInvalTokCondK\,k token、\AmInvalSCond\ 秒，独自探索的智能体为 \AmInvalTokAloneK\,k token、\AmInvalSAlone\ 秒，正确率 \AmInvalAccCond\% 对 \AmInvalAccAlone\%；而在有效或已修复定义上的任务只花 \AmValidTokCondK\,k token。这些使用者找到的修正在每道题里都要重新付一次；把它们作为候选修复经准入接收，就能让这一层对每个修正只付一次。}
